% GENERATED FILE. Edit canonical lessons, then run npm run generate:books.
% canonical-source-hash: fnv1a64:d2a8620e4f3c8cb7
% canonical-lessons: JA-C180-kissaten, JA-C180-shokudou, JA-C180-misoshiru, JA-C180-nigiru, JA-C180-onigiri

\chapter{Cafes, Canteens and Rice Balls}
\label{ch:ja-cafes-canteens-and-rice-balls}
\hlchaptermodality{180}

\begin{chapteropening}
\textbf{By the end of this chapter:} \emph{I can talk about eating out and simple Japanese dishes.}

Complete the last lesson of chapter 180: I can talk about eating out and simple Japanese dishes.
\end{chapteropening}

\section[\texorpdfstring{\ja{きっさてん} (kissaten)}{きっさてん (kissaten)}]{\ja{きっさてん} (kissaten) — cafe; coffee shop}

\label{lesson:JA-C180-kissaten}



\begin{quote}
Before the new one: say the Japanese for to spill, then the Japanese for to warm up.
\end{quote}

\subsection*{You'll want to know: \ja{きっさてん}}
\textbf{\ja{きっさてん}} — \emph{kissaten} — \textquotedblleft{}cafe; coffee shop\textquotedblright{}.

\begin{quote}
\textbf{\ja{ともだちと} \ja{きっさてんで} \ja{コーヒー}\ja{を} \ja{のみます。}} — \emph{tomodachi to kissaten de kōhī o nomimasu} — \textquotedblleft{}I drink coffee at a cafe with a friend.\textquotedblright{}
\end{quote}

\begin{cousinweb}
\textbf{\ja{きっさてん}} (\emph{kissaten}) is Sino-Japanese (kango): three Chinese characters read with their on'yomi, readings borrowed from Chinese. They are \emph{kis}, \textquotedblleft{}consume, drink\textquotedblright{} (Middle Chinese \emph{khek}), \emph{sa}, \textquotedblleft{}tea\textquotedblright{} (\emph{drae}) and \emph{ten}, \textquotedblleft{}shop\textquotedblright{} (\emph{temH}) (KANJIDIC2; Middle Chinese in Baxter's transcription).
\end{cousinweb}

\begin{grammarlens}[title={What you've built}]
One more word for everyday life at A2.
\end{grammarlens}

\subsection*{Your turn}
\begin{itemize}
\raggedright
  \item \emph{Say it:} \emph{kissaten}
  \item \emph{Say it:} \emph{kissaten}, once more
  \item \emph{Say it:} \emph{tomodachi to kissaten de kōhī o nomimasu} — I drink coffee at a cafe with a friend.
  \item \emph{Recall:} say the Japanese for to spill, then the Japanese for to warm up, then say \emph{kissaten} again
\end{itemize}

\begin{tcolorbox}[breakable,colback=teal!4,colframe=teal!35!black,title={Before you move on}]
What does \ja{きっさてん} mean? (\textquotedblleft{}Cafe; coffee shop\textquotedblright{}.) Say it once more.
\end{tcolorbox}
\section[\texorpdfstring{\ja{しょくどう} (shokudō)}{しょくどう (shokudō)}]{\ja{しょくどう} (shokudō) — cafeteria; dining hall, canteen}

\label{lesson:JA-C180-shokudou}



\begin{quote}
Before the new one: say the Japanese for to warm up, then the Japanese for cafe.
\end{quote}

\subsection*{You'll want to know: \ja{しょくどう}}
\textbf{\ja{しょくどう}} — \emph{shokudō} — \textquotedblleft{}cafeteria; dining hall, canteen\textquotedblright{}.

\begin{quote}
\textbf{\ja{がっこうの} \ja{しょくどうで} \ja{ひるごはんを} \ja{たべます。}} — \emph{gakkō no shokudō de hirugohan o tabemasu} — \textquotedblleft{}I eat lunch in the school cafeteria.\textquotedblright{}
\end{quote}

\begin{cousinweb}
\textbf{\ja{しょくどう}} (\emph{shokudō}) is Sino-Japanese (kango): two Chinese characters read with their on'yomi, readings borrowed from Chinese. They are \emph{shoku}, \textquotedblleft{}eat, food\textquotedblright{} (Middle Chinese \emph{zyik}) and \emph{dō}, \textquotedblleft{}hall\textquotedblright{} (\emph{dang}) (KANJIDIC2; Middle Chinese in Baxter's transcription).
\end{cousinweb}

\begin{grammarlens}[title={What you've built}]
One more word for everyday life at A2.
\end{grammarlens}

\subsection*{Your turn}
\begin{itemize}
\raggedright
  \item \emph{Say it:} \emph{shokudō}
  \item \emph{Say it:} \emph{shokudō}, once more
  \item \emph{Say it:} \emph{gakkō no shokudō de hirugohan o tabemasu} — I eat lunch in the school cafeteria.
  \item \emph{Recall:} say the Japanese for to warm up, then the Japanese for cafe, then say \emph{shokudō} again
\end{itemize}

\begin{tcolorbox}[breakable,colback=teal!4,colframe=teal!35!black,title={Before you move on}]
What does \ja{しょくどう} mean? (\textquotedblleft{}Cafeteria; dining hall, canteen\textquotedblright{}.) Say it once more.
\end{tcolorbox}
\section[\texorpdfstring{\ja{みそしる} (misoshiru)}{みそしる (misoshiru)}]{\ja{みそしる} (misoshiru) — miso soup}

\label{lesson:JA-C180-misoshiru}



\begin{quote}
Before the new one: say the Japanese for cafe, then the Japanese for cafeteria.
\end{quote}

\subsection*{You'll want to know: \ja{みそしる}}
\textbf{\ja{みそしる}} — \emph{misoshiru} — \textquotedblleft{}miso soup\textquotedblright{}.

\begin{quote}
\textbf{\ja{あさごはんに} \ja{みそしるを} \ja{のみます。}} — \emph{asagohan ni misoshiru o nomimasu} — \textquotedblleft{}I have miso soup for breakfast.\textquotedblright{}
\end{quote}

\begin{cousinweb}
\textbf{\ja{みそしる}} (\emph{misoshiru}) is built from \emph{miso}, \textquotedblleft{}soybean paste\textquotedblright{}, and native \emph{shiru}, \textquotedblleft{}soup, juice\textquotedblright{}, inherited rather than borrowed (JMdict, KANJIDIC2).
\end{cousinweb}

\begin{grammarlens}[title={What you've built}]
One more word for everyday life at A2.
\end{grammarlens}

\subsection*{Your turn}
\begin{itemize}
\raggedright
  \item \emph{Say it:} \emph{misoshiru}
  \item \emph{Say it:} \emph{misoshiru}, once more
  \item \emph{Say it:} \emph{asagohan ni misoshiru o nomimasu} — I have miso soup for breakfast.
  \item \emph{Recall:} say the Japanese for cafe, then the Japanese for cafeteria, then say \emph{misoshiru} again
\end{itemize}

\begin{tcolorbox}[breakable,colback=teal!4,colframe=teal!35!black,title={Before you move on}]
What does \ja{みそしる} mean? (\textquotedblleft{}Miso soup\textquotedblright{}.) Say it once more.
\end{tcolorbox}
\section[\texorpdfstring{\ja{にぎる} (nigiru)}{にぎる (nigiru)}]{\ja{にぎる} (nigiru) — to grip; to grasp, to shape (rice)}

\label{lesson:JA-C180-nigiru}



\begin{quote}
Before the new one: say the Japanese for cafeteria, then the Japanese for miso soup.
\end{quote}

\subsection*{You'll want to know: \ja{にぎる}}
\textbf{\ja{にぎる}} — \emph{nigiru} — \textquotedblleft{}to grip; to grasp, to shape (rice)\textquotedblright{}.

\begin{quote}
\textbf{\ja{こどもが} \ja{ははの} \ja{てを} \ja{にぎります。}} — \emph{kodomo ga haha no te o nigirimasu} — \textquotedblleft{}The child holds his mother's hand.\textquotedblright{}
\end{quote}

Polite \textbf{\ja{にぎります}}, \emph{nigirimasu}; a godan verb.

\begin{cousinweb}
\textbf{\ja{にぎる}} (\emph{nigiru}) is a native Japanese word (wago), inherited rather than borrowed from Chinese; its kanji was attached later as a kun'yomi (JMdict, KANJIDIC2).
\end{cousinweb}

\begin{grammarlens}[title={What you've built}]
One more word for everyday life at A2.
\end{grammarlens}

\subsection*{Your turn}
\begin{itemize}
\raggedright
  \item \emph{Say it:} \emph{nigiru}
  \item \emph{Say it:} \emph{nigiru}, once more
  \item \emph{Say it:} \emph{kodomo ga haha no te o nigirimasu} — The child holds his mother's hand.
  \item \emph{Recall:} say the Japanese for cafeteria, then the Japanese for miso soup, then say \emph{nigiru} again
\end{itemize}

\begin{tcolorbox}[breakable,colback=teal!4,colframe=teal!35!black,title={Before you move on}]
What does \ja{にぎる} mean? (\textquotedblleft{}To grip; to grasp, to shape (rice)\textquotedblright{}.) Say it once more.
\end{tcolorbox}
\section[\texorpdfstring{\ja{おにぎり} (onigiri)}{おにぎり (onigiri)}]{\ja{おにぎり} (onigiri) — rice ball}

\label{lesson:JA-C180-onigiri}



\begin{quote}
Before the new one: say the Japanese for miso soup, then the Japanese for to grip.
\end{quote}

\subsection*{You'll want to know: \ja{おにぎり}}
\textbf{\ja{おにぎり}} — \emph{onigiri} — \textquotedblleft{}rice ball\textquotedblright{}.

\begin{quote}
\textbf{\ja{ははは} \ja{まいあさ} \ja{おにぎりを} \ja{にぎります。}} — \emph{haha wa maiasa onigiri o nigirimasu} — \textquotedblleft{}My mother makes rice balls every morning.\textquotedblright{}
\end{quote}

\begin{cousinweb}
\textbf{\ja{おにぎり}} (\emph{onigiri}) is built from polite \emph{o-} and native \emph{nigiri}, \textquotedblleft{}a grip, something shaped in the hand\textquotedblright{}, which comes from \emph{nigiru}, \textquotedblleft{}to grip\textquotedblright{} (JMdict, KANJIDIC2).
\end{cousinweb}

\begin{grammarlens}[title={What you've built}]
One more word for everyday life at A2.
\end{grammarlens}

\subsection*{Your turn}
\begin{itemize}
\raggedright
  \item \emph{Say it:} \emph{onigiri}
  \item \emph{Say it:} \emph{onigiri}, once more
  \item \emph{Say it:} \emph{haha wa maiasa onigiri o nigirimasu} — My mother makes rice balls every morning.
  \item \emph{Recall:} say the Japanese for miso soup, then the Japanese for to grip, then say \emph{onigiri} again
\end{itemize}

\begin{tcolorbox}[breakable,colback=teal!4,colframe=teal!35!black,title={Before you move on}]
What does \ja{おにぎり} mean? (\textquotedblleft{}Rice ball\textquotedblright{}.) Say it once more.
\end{tcolorbox}
